\documentclass[11pt]{article}
\usepackage[T1]{fontenc}
\usepackage[utf8]{inputenc}
\usepackage{lmodern}
\usepackage{amsmath,amssymb,amsthm}
\usepackage[margin=2.6cm]{geometry}
\usepackage{booktabs}
\usepackage{enumitem}
\usepackage{xcolor}
\usepackage[colorlinks=true,linkcolor=blue!50!black,citecolor=blue!50!black,urlcolor=blue!50!black]{hyperref}

\setlist{itemsep=2pt,topsep=4pt}
\setlength{\emergencystretch}{2em}
\allowdisplaybreaks[1]

\newtheorem{theorem}{Theorem}[section]
\newtheorem{lemma}[theorem]{Lemma}
\newtheorem{proposition}[theorem]{Proposition}
\newtheorem{corollary}[theorem]{Corollary}
\newtheorem{conjecture}[theorem]{Conjecture}
\theoremstyle{definition}
\newtheorem{definition}[theorem]{Definition}
\theoremstyle{remark}
\newtheorem{remark}[theorem]{Remark}

\newcommand{\Hm}{\mathrm{H}}
\newcommand{\D}{D}
\newcommand{\tht}{\theta}
\newcommand{\tth}{\tilde\theta}
\newcommand{\Bad}{\mathrm{Bad}}

\title{\textbf{A proof of Liu's conjecture on the harmonic index\\ and the diameter of a graph}\\[4pt]
\large The inequality $\Hm(G)\ge \tfrac12\D(G)+\tfrac13$ via a vertex-deletion calculus}
\author{Verifiable Mathematical Reasoning Corpus --- research manuscript}
\date{20 November 2026}

\begin{document}
\maketitle

\begin{abstract}
The harmonic index of a graph $G$ is $\Hm(G)=\sum_{uv\in E(G)}2/(d_u+d_v)$. In 2013 J.~Liu
conjectured that every connected graph of order $n\ge4$ and diameter $\D$ satisfies
$\Hm(G)\ge \D+\tfrac56-\tfrac n2$ and $\Hm(G)\ge\bigl(\tfrac12+\tfrac1{3(n-1)}\bigr)\D$, the path
being extremal. The conjecture was known for trees, unicyclic graphs and quasi-tree graphs. We
prove a statement that does not involve the order at all: for every connected graph with at
least two vertices,
\[
 \Hm(G)\ \ge\ \tfrac12\,\D(G)+\tfrac13 ,
\]
with equality if and only if $G$ is a path with at least three vertices. Both inequalities of
Liu's conjecture, together with the characterization of equality, follow at once. The proof
is an induction on the order which deletes one vertex at a time. It rests on an exact formula
for the sum of the changes $\Hm(G)-\Hm(G-z)$ over a set of vertices, on a lemma which removes
every part of the graph hanging off a diametral path, and on a \emph{block lemma}: in a
$2$-connected block with two distinguished vertices, the deletions of the remaining vertices
decrease the harmonic index on average. The manuscript has not been refereed; every identity
and inequality of the proof has additionally been checked in exact arithmetic on all connected
graphs of small order, and the main inequality on all $1{,}006{,}700{,}565$ connected graphs
of order $11$.
\end{abstract}

\section{Introduction}\label{sec:intro}

\subsection{The conjecture}
All graphs are finite and simple. For a connected graph $G$ we write $n=|V(G)|$, $d_v$ for the
degree of $v$, $d(x,y)$ for the distance and $\D(G)=\max_{x,y}d(x,y)$ for the diameter. The
\emph{harmonic index}
\[
 \Hm(G)=\sum_{uv\in E(G)}\frac{2}{d_u+d_v}
\]
goes back to Fajtlowicz's Graffiti conjectures \cite{Faj} and was studied by Favaron, Mah\'eo
and Sacl\'e \cite{FMS}; see the survey \cite{AZG} for the extensive literature. Zhong
\cite{Zhong} proved that the star minimizes $\Hm$ among connected graphs:
\begin{equation}\label{eq:zhong}
 \Hm(G)\ \ge\ \frac{2(n-1)}{n}\qquad\text{for every connected graph $G$ of order $n\ge2$}
\end{equation}
(for $n=2$ the graph is $K_2$ and both sides equal $1$).
Relating $\Hm$ to the diameter, Liu \cite{Liu} established the two inequalities below for
trees and proposed them for all connected graphs.

\begin{conjecture}[Liu \cite{Liu}, 2013]\label{conj:liu}
Let $G$ be a connected graph of order $n\ge4$ with diameter $\D(G)$. Then
\[
 \Hm(G)\ \ge\ \D(G)+\frac56-\frac n2
 \qquad\text{and}\qquad
 \Hm(G)\ \ge\ \Bigl(\frac12+\frac1{3(n-1)}\Bigr)\D(G) ,
\]
and the bounds are attained by the path $P_n$.
\end{conjecture}

Jerline and Michaelraj proved the conjecture for unicyclic graphs, with sharper constants
\cite{JM1,JM2}, and Abdolghafourian and Iranmanesh for quasi-tree graphs \cite{AI}. We are
not aware of a proof for general graphs; the scope of our literature search is described in
Section~\ref{sec:status}.

\subsection{Results}
For the path, $\Hm(P_n)=\frac n2-\frac16=\frac12\D(P_n)+\frac13$ when $n\ge3$. Our main
result says that no connected graph has a smaller harmonic index relative to its diameter.

\begin{theorem}\label{thm:main}
Let $G$ be a connected graph with at least two vertices. Then
\[
 \Hm(G)\ \ge\ \frac12\,\D(G)+\frac13 ,
\]
with equality if and only if $G$ is a path with at least three vertices.
\end{theorem}

\begin{corollary}[Liu's conjecture]\label{cor:liu}
Let $G$ be a connected graph of order $n\ge3$ with diameter $\D$. Then
\[
 \Hm(G)\ \ge\ \D+\frac56-\frac n2
 \qquad\text{and}\qquad
 \Hm(G)\ \ge\ \Bigl(\frac12+\frac1{3(n-1)}\Bigr)\D ,
\]
and in each inequality equality holds if and only if $G\cong P_n$.
\end{corollary}

Theorem~\ref{thm:main} is stronger than Conjecture~\ref{conj:liu}: it does not involve $n$,
and for a graph whose diameter is small compared with its order the two inequalities of the
conjecture are far weaker (the first one is void as soon as $\D\le \frac n2-\frac56$).

\subsection{Method}
The inequality of Theorem~\ref{thm:main} was found, and is proved, by looking at what happens
to $\Hm$ when a single vertex is deleted. Three ingredients are new.
\begin{itemize}
\item \emph{An exact deletion formula} (Lemma~\ref{lem:deletion}): for every set $Z$ of
 vertices, $\sum_{z\in Z}[\Hm(G)-\Hm(G-z)]$ is a sum over the edges of $G$ of explicit terms
 depending only on the two end degrees and on the number of neighbours in $Z$. For $Z=V(G)$
 it gives the identity of Remark~\ref{rem:identity}.
\item \emph{Removal of hanging parts} (Lemma~\ref{lem:hanging}): if a component $X$ of $G-v$
 does not exhaust $G-v$, then $\Hm(G)>\Hm(G-X)$. Hence everything that hangs off a fixed
 diametral path can be deleted, and the graph becomes a chain of blocks along that path.
\item \emph{The block lemma} (Proposition~\ref{prop:block}): in a $2$-connected block of such
 a chain, the sum of $\Hm(G)-\Hm(G-z)$ over all vertices $z$ of the block other than its two
 ``poles'' is positive. Consequently some vertex which is neither a cut vertex nor an end of
 the diametral path can be deleted with a strict decrease of $\Hm$ and without decreasing the
 diameter.
\end{itemize}
An induction on the order then gives Theorem~\ref{thm:main}. No single local rule for choosing
the deleted vertex works in general (Remark~\ref{rem:rules}); the averaging over a block is
what makes the argument uniform. The calculus grew out of our earlier study of the harmonic
form of Fajtlowicz's conjecture on the Randi\'c index and the average distance \cite{VMRC},
where an edge-deletion identity plays a similar role.

\section{Notation and the deletion formula}\label{sec:deletion}

For an edge $uw$ we write $S_{uw}=d_u+d_w$. If $n\ge3$ and $G$ is connected then
$S_{uw}\ge3$ for every edge. We use the two functions
\[
 \tht(S)=\frac{2}{S(S-1)}=\frac2{S-1}-\frac2S ,\qquad
 \tth(S)=\frac{2}{S(S-2)}=\frac12\Bigl(\frac2{S-2}-\frac2S\Bigr)\qquad(S\ge3).
\]
Both are positive and decreasing, $\tht(S)<\tth(S)$, and $(S-2)\,\tth(S)=2/S$. For a graph
$G'$ which is not connected, $\Hm(G')$ is defined by the same sum over its edges; isolated
vertices contribute nothing.

Let $Z\subseteq V(G)$ and let $uw$ be an edge. We put
\begin{align*}
 j_Z(uw)&=|\{u,w\}\cap Z| ,\\
 k_Z(uw)&=\text{the number of $z\in Z\setminus\{u,w\}$ adjacent to exactly one of $u,w$},\\
 t_Z(uw)&=\text{the number of $z\in Z$ adjacent to both $u$ and $w$},\\
 f_Z(uw)&=|N(u)\setminus(Z\cup\{w\})|+|N(w)\setminus(Z\cup\{u\})| .
\end{align*}
Counting the neighbours of $u$ other than $w$ and of $w$ other than $u$ gives
\begin{equation}\label{eq:count}
 k_Z(uw)+2\,t_Z(uw)=S_{uw}-2-f_Z(uw).
\end{equation}

\begin{lemma}[deletion formula]\label{lem:deletion}
Let $G$ be a connected graph of order $n\ge3$ and $Z\subseteq V(G)$. Then
\begin{align*}
 \sum_{z\in Z}\bigl[\Hm(G)-\Hm(G-z)\bigr]&=\sum_{uw\in E(G)}\omega_Z(uw),\\
 \omega_Z(uw)&=j_Z(uw)\frac{2}{S_{uw}}-k_Z(uw)\,\tht(S_{uw})-2\,t_Z(uw)\,\tth(S_{uw}).
\end{align*}
Moreover, for every edge $uw$, with $S=S_{uw}$ and $f=f_Z(uw)$,
\begin{equation}\label{eq:omega}
 \omega_Z(uw)\ \ge\ \tth(S)\bigl[(j_Z(uw)-1)(S-2)+f\bigr],
\end{equation}
that is, $\omega_Z\ge \frac2S+f\,\tth(S)$ if both ends are in $Z$, $\omega_Z\ge f\,\tth(S)$ if
exactly one end is in $Z$, and $\omega_Z\ge-(S-2-f)\,\tth(S)$ if no end is in $Z$. If
$t_Z(uw)=0$ then $\omega_Z(uw)=\tht(S)\bigl[(j_Z(uw)-1)(S-1)+1+f\bigr]$.
\end{lemma}
\begin{proof}
Fix $z$. Deleting $z$ removes the edges at $z$, of total weight $\sum_{u\sim z}2/S_{zu}$. An
edge $uw$ not incident with $z$ survives, and its degree sum drops by the number
$c\in\{0,1,2\}$ of its ends adjacent to $z$; its weight increases by $2/(S-1)-2/S=\tht(S)$ if
$c=1$ and by $2/(S-2)-2/S=2\tth(S)$ if $c=2$. Summing over $z\in Z$ and collecting the
contributions of each edge gives the formula. By \eqref{eq:count} and $\tht\le\tth$,
$k\tht+2t\tth\le(S-2-f)\tth$, and $2/S=(S-2)\tth(S)$; this is \eqref{eq:omega}. If $t=0$ then
$k=S-2-f$ and $\omega=j\cdot\frac2S-(S-2-f)\tht(S)$, which is the last formula because
$\frac2S=(S-1)\tht(S)$.
\end{proof}

\begin{remark}[a vertex-deletion identity]\label{rem:identity}
For $Z=V(G)$ we have $j=2$, $f=0$ and $k=S-2-2t$, where $t=t_{uw}$ is the number of common
neighbours of $u$ and $w$, so
\[
 \sum_{z\in V(G)}\bigl[\Hm(G)-\Hm(G-z)\bigr]
 =\sum_{uw\in E(G)}\Bigl[\frac{2}{S_{uw}-1}-\frac{4\,t_{uw}}{S_{uw}(S_{uw}-1)(S_{uw}-2)}\Bigr]
 \ \ge\ \Hm(G),
\]
the inequality being \eqref{eq:omega}. In particular every connected graph with at least
three vertices has a vertex whose deletion strictly decreases the harmonic index. The difficulty in using this for
Theorem~\ref{thm:main} is that the deleted vertex must not destroy the diameter.
\end{remark}

\section{Hanging parts and the chain structure}\label{sec:hanging}

\begin{lemma}[hanging parts]\label{lem:hanging}
Let $G$ be a connected graph, $v\in V(G)$, and let $X$ be the vertex set of a component of
$G-v$ such that $v$ has a neighbour outside $X$. Then $\Hm(G)>\Hm(G-X)$.
\end{lemma}
\begin{proof}
Let $e=|N(v)\cap X|\ge1$ and $x=d_v-e\ge1$. In $G-X$ only the degree of $v$ changes, from
$e+x$ to $x$. Hence $\Hm(G)-\Hm(G-X)=A-B$ with
\[
 A=\sum_{\substack{uw\in E(G)\\ \{u,w\}\cap X\ne\emptyset}}\frac2{S_{uw}},\qquad
 B=\sum_{u\in N(v)\setminus X}\Bigl[\frac{2}{x+d_u}-\frac2{x+e+d_u}\Bigr]
 =\sum_{u\in N(v)\setminus X}\frac{2e}{(x+d_u)(x+e+d_u)} .
\]
Let $J=G[X\cup\{v\}]$. It is connected, has $|X|+1\ge e+1$ vertices, the degree of $v$ in $J$
is $e$ and all other degrees are as in $G$. Therefore
\[
 \Hm(J)=A+\sum_{y\in N(v)\cap X}\Bigl[\frac2{e+d_y}-\frac2{e+x+d_y}\Bigr]
 =A+\sum_{y\in N(v)\cap X}\frac{2x}{(e+d_y)(e+x+d_y)} .
\]
By \eqref{eq:zhong}, $\Hm(J)\ge 2|X|/(|X|+1)\ge 2e/(e+1)$. Using $d_y\ge1$ and $d_u\ge1$,
\begin{align*}
 A-B\ &\ge\ \frac{2e}{e+1}-\frac{2xe}{(e+1)(e+x+1)}-\frac{2xe}{(x+1)(e+x+1)}\\
 &=\frac{2e\bigl[(x+1)(e+x+1)-x(x+1)-x(e+1)\bigr]}{(e+1)(x+1)(e+x+1)} .
\end{align*}
Since $(x+1)(e+x+1)=(x+1)(e+1)+x(x+1)$, the bracket equals $e+1$, and
$A-B\ge 2e/((x+1)(e+x+1))>0$.
\end{proof}

Now fix a connected graph $G$ with $n\ge3$, two vertices $p,q$ with $d(p,q)=\D(G)$, and a
shortest path $P=v_0v_1\cdots v_\D$ from $p=v_0$ to $q=v_\D$. We say that $G$ has \emph{no
hanging part} (with respect to $P$) if
\begin{equation}\label{eq:H}
 \text{for every vertex $v$, every component of $G-v$ contains a vertex of $P$.}\tag{H}
\end{equation}
Recall that a \emph{block} is a maximal connected subgraph without a cut vertex; a block is a
$K_2$ or is $2$-connected with at least three vertices; a vertex lies in more than one block
if and only if it is a cut vertex; a block containing no cut vertex of the connected graph $G$
is $G$ itself; and if two blocks share a vertex $a$, their remaining vertices lie in different
components of $G-a$. Note that $p\ne q$ because $n\ge3$.

\begin{lemma}[chain structure]\label{lem:chain}
Assume \eqref{eq:H}. Then:
\begin{enumerate}[label=\textup{(\roman*)}]
\item Every cut vertex $v$ of $G$ is an internal vertex $v_i$ of $P$, and $G-v_i$ has exactly
 two components, one containing $v_0,\dots,v_{i-1}$ and the other $v_{i+1},\dots,v_\D$. Every
 vertex of degree one is $p$ or $q$.
\item Let $\Gamma$ be a block with at least three vertices, and let $\Pi$ be the set of all
 cut vertices of $G$ in $\Gamma$ together with those of $p,q$ that lie in $\Gamma$. Then
 $|\Pi|=2$.
\item With $\Gamma,\Pi$ as in \textup{(ii)} put $Z=V(\Gamma)\setminus\Pi$. Then $Z\ne\emptyset$,
 no vertex of $Z$ is a cut vertex or one of $p,q$, and every neighbour of a vertex of $Z$ lies
 in $\Gamma$.
\item For $a\in\Pi$ let $N_{\rm out}(a)=N(a)\setminus V(\Gamma)$. Then either
 $N_{\rm out}(a)=\emptyset$, or $N_{\rm out}(a)$ is a single vertex, or $|N_{\rm out}(a)|\ge2$
 and every vertex of $N_{\rm out}(a)$ has degree at least two.
\end{enumerate}
\end{lemma}
\begin{proof}
(i) $P-v$ has at most two components, and exactly two only if $v$ is an internal vertex of
$P$. By \eqref{eq:H} every component of $G-v$ contains a component of $P-v$; so a cut vertex is
an internal vertex $v_i$, and $G-v_i$ has exactly two components, which contain the two parts
of $P-v_i$. If $d_y=1$ and $y\notin\{p,q\}$, let $v$ be the neighbour of $y$; then $v\ne y$
has a further neighbour (as $n\ge3$), and $\{y\}$ is a component of $G-v$; by \eqref{eq:H},
$y\in V(P)$, and an internal vertex of $P$ has degree at least two, a contradiction.

(ii) $\Gamma$ is $2$-connected, so $\Gamma-v$ is connected for every $v\in V(\Gamma)$. If
$\Gamma$ contains no cut vertex of $G$, then $\Gamma=G$ and $\Pi=\{p,q\}$. Suppose $\Gamma$
contains three cut vertices $v_i,v_j,v_k$ with $i<j<k$. Then $\Gamma-v_j$ contains $v_i$ and
$v_k$, which by (i) lie in different components of $G-v_j$; a contradiction. If $\Gamma$
contains two cut vertices $v_i,v_j$ with $i<j$, the same argument shows $p\notin V(\Gamma)$
(otherwise $\Gamma-v_i$ would contain $p$ and $v_j$) and $q\notin V(\Gamma)$ (otherwise
$\Gamma-v_j$ would contain $v_i$ and $q$); so $\Pi=\{v_i,v_j\}$. Finally let $\Gamma$ contain
exactly one cut vertex $v_i$, and let $C$ be the component of $G-v_i$ containing
$\Gamma-v_i$. If $C\ne V(\Gamma)\setminus\{v_i\}$, some edge joins a vertex $g$ of
$\Gamma-v_i$ to a vertex of $C$ outside $\Gamma$; this edge lies in another block, so $g$ is a
second cut vertex in $\Gamma$, a contradiction. Thus $V(\Gamma)\setminus\{v_i\}$ is a
component of $G-v_i$, and by (i) it contains exactly one of $p,q$. In all cases
$|\Pi|=2$.

(iii) $|V(\Gamma)|\ge3$ gives $Z\ne\emptyset$. A vertex of $Z$ is not a cut vertex, so it lies
in no other block and all its neighbours are in $\Gamma$; it is not $p$ or $q$ by the
definition of $\Pi$.

(iv) If $a\in\{p,q\}$, then $a$ is not a cut vertex by (i) and $N_{\rm out}(a)=\emptyset$.
Let $a$ be a cut vertex. By (i), $G-a$ has two components; $\Gamma-a$ lies in one of them,
and the vertices other than $a$ of any other block through $a$ lie in a different component,
hence in the second one. Two such blocks would be separated by $a$ from each other, so there
is exactly one further block $\Gamma'$ through $a$, and $N_{\rm out}(a)=N(a)\cap V(\Gamma')$.
If $\Gamma'=K_2$ this is one vertex. Otherwise $\Gamma'$ is $2$-connected, $a$ has at least
two neighbours in it, and each of them has degree at least two.
\end{proof}

The two vertices of $\Pi$ are called the \emph{poles} of $\Gamma$.

\section{The block lemma}\label{sec:block}

Throughout this section $G$ is a connected graph of order $n\ge3$, $\Gamma$ is a block of $G$
with at least three vertices, $\Pi=\{a_1,a_2\}\subseteq V(\Gamma)$ and
$Z=V(\Gamma)\setminus\Pi$, and we assume:
\begin{enumerate}[label=\textup{(B\arabic*)}]
\item every neighbour of a vertex of $Z$ lies in $\Gamma$;
\item for each $a\in\Pi$ the set $N_{\rm out}(a)=N(a)\setminus V(\Gamma)$ is empty, or a single
 vertex, or all its vertices have degree at least two.
\end{enumerate}
By Lemma~\ref{lem:chain} these hold for every block with at least three vertices of a graph
without hanging parts. We use the following notation. For $a\in\Pi$, $a'$ is the other pole,
\[
 n_a=|N(a)\cap Z|,\qquad x_a=|N_{\rm out}(a)|,\qquad \varepsilon=\begin{cases}1&a_1a_2\in E(G)\\0&\text{otherwise,}\end{cases}
 \qquad O_a=\sum_{u'\in N_{\rm out}(a)}\tht(d_a+d_{u'}),
\]
so that $d_a=n_a+x_a+\varepsilon$. For $u\in Z$, $\pi_u=|N(u)\cap\Pi|\in\{0,1,2\}$, and
$\Delta=\max_{u\in Z}d_u$. Since $\Gamma$ is $2$-connected, every vertex of $\Gamma$ has at
least two neighbours in $\Gamma$; thus $d_u\ge2$ for $u\in Z$, $d_a\ge2$, and
$n_a+\varepsilon\ge2$ for $a\in\Pi$. An edge with both ends in $Z$ is called a
\emph{$ZZ$-edge}.

\begin{proposition}[block lemma]\label{prop:block}
Under \textup{(B1)} and \textup{(B2)},
\[
 \Sigma_\Gamma:=\sum_{z\in Z}\bigl[\Hm(G)-\Hm(G-z)\bigr]\ >\ 0 .
\]
In particular $\Hm(G-z)<\Hm(G)$ for some $z\in Z$.
\end{proposition}

The proof occupies the rest of the section.

\begin{lemma}\label{lem:lower}
\[
 \Sigma_\Gamma\ \ge\ \sum_{uw\ ZZ\text{-edge}}\Bigl[\frac{2}{S_{uw}}+(\pi_u+\pi_w)\tth(S_{uw})\Bigr]
 +\sum_{a\in\Pi}\ \sum_{u\in N(a)\cap Z}\Bigl[\omega_Z(au)-O_a-\varepsilon\,\tth(d_{a_1}+d_{a_2})\Bigr],
\]
and $\omega_Z(au)\ge(x_a+\varepsilon+\pi_u-1)\,\tth(d_a+d_u)$ for $a\in\Pi$, $u\in N(a)\cap Z$.
\end{lemma}
\begin{proof}
We apply Lemma~\ref{lem:deletion} to the set $Z$ and go through the edges of $G$.

\emph{$ZZ$-edges.} By (B1) the neighbours of $u\in Z$ outside $Z$ are its $\pi_u$ poles, so
$f_Z(uw)=\pi_u+\pi_w$ and \eqref{eq:omega} gives $\omega_Z(uw)\ge 2/S_{uw}+(\pi_u+\pi_w)\tth(S_{uw})$.

\emph{Edges $au$ with $a\in\Pi$, $u\in Z$.} The neighbours of $a$ outside $Z$ and different from
$u$ are its $x_a$ outside neighbours and, if $\varepsilon=1$, the other pole; those of $u$
different from $a$ are $\pi_u-1$ in number. So $f_Z(au)=x_a+\varepsilon+\pi_u-1$ and
\eqref{eq:omega} gives the stated bound for $\omega_Z(au)$.

\emph{The edge $a_1a_2$} (if $\varepsilon=1$). Here $j_Z=0$ and, by \eqref{eq:count},
$S-2-f_Z=k_Z+2t_Z=n_{a_1}+n_{a_2}$, so
$\omega_Z(a_1a_2)\ge-(n_{a_1}+n_{a_2})\tth(d_{a_1}+d_{a_2})$ by \eqref{eq:omega}.

\emph{Edges $au'$ with $a\in\Pi$, $u'\in N_{\rm out}(a)$.} By (B1), $u'$ has no neighbour in
$Z$. Hence $j_Z=0$, $t_Z=0$, $k_Z=n_a$ and $\omega_Z(au')=-n_a\tht(d_a+d_{u'})$.

\emph{All other edges} have no end in $Z$ and, by (B1), no end adjacent to a vertex of $Z$;
for them $\omega_Z=0$.

Adding up, and writing $n_{a_1}+n_{a_2}$ and $n_aO_a$ as sums over the pairs $(a,u)$ with
$u\in N(a)\cap Z$, we obtain the lemma.
\end{proof}

The next lemma is a statement about numbers; in it $n$ denotes an arbitrary positive integer
(it will be applied with $n=n_a$), not the order of $G$. Recall also that $\Delta$ is the
maximum degree over $Z$ only, not the maximum degree of $G$.

\begin{lemma}[pole inequality]\label{lem:pole}
Let $\Delta\ge3$, $n\ge1$, $x\ge0$ be integers, $\varepsilon\in\{0,1\}$ with $n+\varepsilon\ge2$,
$d=n+x+\varepsilon$, and let $d'\ge2$ and $O\ge0$ be real numbers. Put
\[
 \beta=n\bigl[O+\varepsilon\,\tth(d+d')-(x+\varepsilon)\,\tth(d+\Delta)\bigr].
\]
\begin{enumerate}[label=\textup{(\alph*)}]
\item If $O\le x\,\tht(d+2)$, then $\beta<\frac12-\frac1\Delta$.
\item If $x=1$ and $O=\tht(d+1)$, then $\beta\le\frac12-\frac1\Delta$, with equality only for
 $\varepsilon=0$, $n=2$, $\Delta=3$.
\end{enumerate}
\end{lemma}
\begin{proof}
(a) Let $m=x+\varepsilon$. If $m=0$ then $\beta=0$. Let $m\ge1$, so $d=n+m\ge2$. Since
$\tht\le\tth$, $\tth$ is decreasing and $d'\ge2$,
\begin{align*}
 \beta&\le nm\bigl[\tth(d+2)-\tth(d+\Delta)\bigr]\le\frac{d^2}{4}\Bigl[\frac{2}{d(d+2)}-\frac{2}{(d+\Delta)(d+\Delta-2)}\Bigr]\\
 &=\frac12-\frac1{d+2}-\frac{d^2}{2(d+\Delta)(d+\Delta-2)} ,
\end{align*}
where we used $nm\le(n+m)^2/4$. It remains to show
\begin{equation}\label{eq:II}
 \frac1\Delta<\frac1{d+2}+\frac{d^2}{2(d+\Delta)(d+\Delta-2)}\qquad(d\ge2,\ \Delta\ge3).
\end{equation}
If $\Delta\ge d+2$ this is clear. Let $3\le\Delta\le d+1$. Because
$(d+\Delta)(d+\Delta-2)<(d+\Delta-1)^2$, it suffices that
$g(\Delta):=d^2\Delta(d+2)-2(d+\Delta-1)^2(d+2-\Delta)\ge0$ on the real interval $[3,d+1]$.
Now $g(3)=(d+2)(d^2-2d+4)>0$, and
\[
 g'(\Delta)=d^2(d+2)+2(d+\Delta-1)(3\Delta-d-5).
\]
For $3\le\Delta\le d+1$ we have $0<d+\Delta-1\le2d$ and $3\Delta-d-5\ge-(d-4)$. If $d\le4$
then $g'>0$; if $d\ge5$ then $g'(\Delta)\ge d^2(d+2)-4d(d-4)=d(d^2-2d+16)>0$. Hence $g$ is
increasing and positive on $[3,d+1]$, which proves \eqref{eq:II}.

(b) \emph{Case $\varepsilon=0$.} Then $d=n+1$, $n\ge2$ and
\[
 \beta=\frac{2n}{(n+1)(n+2)}-\frac{2n}{(n+\Delta+1)(n+\Delta-1)} .
\]
Note that $\frac{2n}{(n+1)(n+2)}\le\frac13$ for $n\ge2$ (equivalently $(n-1)(n-2)\ge0$) and
$\frac{2n}{(n+1)(n+2)}\le\frac3{10}$ for $n\ge3$ (equivalently $(3n-2)(n-3)\ge0$). So
$\beta<\frac13\le\frac12-\frac1\Delta$ for $\Delta\ge6$. For $\Delta=5$:
$\beta<\frac3{10}$ if $n\ge3$, and $\beta=\frac14$ if $n=2$. For $\Delta=4$:
$\beta=\frac{2n(5n+13)}{(n+1)(n+2)(n+3)(n+5)}$, which is $\frac{92}{420}$ and
$\frac{168}{960}$ for $n=2,3$, and for $n\ge4$ it is less than $\frac14$ because
$(n+2)(n+3)(n+1)(n+5)\ge42(n^2+6n+5)>8n(5n+13)$. For $\Delta=3$:
$\beta=\frac{6n}{(n+1)(n+2)(n+4)}$ and
\[
 \frac16-\beta=\frac{(n-2)(n^2+9n-4)}{6(n+1)(n+2)(n+4)}\ \ge0 ,
\]
with equality only for $n=2$.

\emph{Case $\varepsilon=1$.} Then $d=n+2$, $n\ge1$, and using $d'\ge2$,
\begin{align*}
 \beta&\le c(n)-\frac{4n}{(n+\Delta)(n+\Delta+2)},\\
 c(n)&=\frac{2n}{(n+2)(n+3)}+\frac{2n}{(n+2)(n+4)}=\frac{2n(2n+7)}{(n+2)(n+3)(n+4)} .
\end{align*}
One checks that $c(n+1)<c(n)$ if and only if $n^2+2n-9>0$, i.e.\ $n\ge3$; hence
$c(n)\le c(3)=\frac{13}{35}$ for all $n$, and $c$ is decreasing for $n\ge3$. For
$\Delta\ge8$ we get $\beta<\frac{13}{35}<\frac38\le\frac12-\frac1\Delta$. For
$\Delta\in\{4,5,6,7\}$ let $N_\Delta=10,7,6,4$ respectively; then
$c(N_\Delta)=\frac{540}{2184},\frac{294}{990},\frac{228}{720},\frac{5}{14}$ is at most
$\frac12-\frac1\Delta=\frac14,\frac3{10},\frac13,\frac5{14}$, so $\beta<\frac12-\frac1\Delta$
for $n\ge N_\Delta$; for the finitely many $n<N_\Delta$ the numbers
$c(n)-\frac{4n}{(n+\Delta)(n+\Delta+2)}$ are listed, as exact fractions, in
Table~\ref{tab:pole}; each of them is smaller than $\frac12-\frac1\Delta$ (the largest entry
of each row is $\frac15$, $\frac{151}{630}$, $\frac4{15}$, $\frac{283}{990}$, to be compared
with $\frac14$, $\frac3{10}$, $\frac13$, $\frac5{14}$). Finally, for $\Delta=3$,
\[
 c(n)-\frac{4n}{(n+3)(n+5)}=\frac{2n(5n+19)}{(n+2)(n+3)(n+4)(n+5)}<\frac16 :
\]
indeed $12n(5n+19)<(n+2)(n+3)(n+4)(n+5)$ holds for $n=1,2$ ($288<360$, $696<840$), for
$3\le n\le17$ because $(n+4)(n+5)\ge56$ and $56(n^2+5n+6)\ge60n^2+228n$, and for $n\ge18$
because $(n+4)(n+5)\ge506$.
\end{proof}

\begin{table}[h]
\centering
\renewcommand{\arraystretch}{1.35}
\begin{tabular}{@{}c|ccccccccc|c@{}}
\toprule
$\Delta\backslash n$ & 1&2&3&4&5&6&7&8&9& $\frac12-\frac1\Delta$\\
\midrule
4 & $\frac{13}{70}$ & $\frac{1}{5}$ & $\frac{19}{105}$ & $\frac{11}{70}$ & $\frac{125}{924}$ & $\frac{7}{60}$ & $\frac{217}{2145}$ & $\frac{34}{385}$ & $\frac{111}{1430}$ &$\frac14$\\
5 & $\frac{13}{60}$ & $\frac{151}{630}$ & $\frac{31}{140}$ & $\frac{271}{1386}$ & $\frac{43}{252}$ & $\frac{1277}{8580}$ & & & &$\frac3{10}$\\
6 & $\frac{149}{630}$ & $\frac{4}{15}$ & $\frac{289}{1155}$ & $\frac{47}{210}$ & $\frac{7115}{36036}$ & & & & &$\frac13$\\
7 & $\frac{1}{4}$ & $\frac{283}{990}$ & $\frac{19}{70}$ & & & & & & &$\frac5{14}$\\
\bottomrule
\end{tabular}
\caption{The exact values of $c(n)-\frac{4n}{(n+\Delta)(n+\Delta+2)}$ for $n<N_\Delta$ (proof
of Lemma~\ref{lem:pole}(b), case $\varepsilon=1$).}\label{tab:pole}
\end{table}

\begin{proof}[Proof of Proposition~\ref{prop:block} when $\Delta\ge3$]
Let $a\in\Pi$ and $u\in N(a)\cap Z$. Then $\pi_u\ge1$ and $d_u\le\Delta$, so by
Lemma~\ref{lem:lower},
\[
 \omega_Z(au)-O_a-\varepsilon\tth(d_{a_1}+d_{a_2})
 \ \ge\ -\bigl[O_a+\varepsilon\tth(d_a+d_{a'})-(x_a+\varepsilon)\tth(d_a+\Delta)\bigr].
\]
Summing over the $n_a$ vertices $u$ shows that the contribution of the pole $a$ in
Lemma~\ref{lem:lower} is at least $-\Bad_a$, where
\[
 \Bad_a=n_a\max\bigl\{0,\ O_a+\varepsilon\tth(d_a+d_{a'})-(x_a+\varepsilon)\tth(d_a+\Delta)\bigr\}.
\]
We claim $\Bad_a\le\frac12-\frac1\Delta$. This is Lemma~\ref{lem:pole} with $n=n_a$, $x=x_a$,
$d=d_a$, $d'=d_{a'}\ge2$, $O=O_a$: if $x_a=0$, or if all vertices of $N_{\rm out}(a)$ have
degree at least two, then $O_a\le x_a\tht(d_a+2)$ and part (a) applies; otherwise, by (B2),
$N_{\rm out}(a)$ is a single vertex of degree one, $x_a=1$, $O_a=\tht(d_a+1)$, and part (b)
applies. Therefore
\begin{equation}\label{eq:final}
 \Sigma_\Gamma\ \ge\ \sum_{uw\ ZZ\text{-edge}}\frac2{S_{uw}}+E-\Bad_{a_1}-\Bad_{a_2}
 \ \ge\ \sum_{uw\ ZZ\text{-edge}}\frac2{S_{uw}}+E-\Bigl(1-\frac2\Delta\Bigr),
\end{equation}
where $E=\sum_{ZZ}(\pi_u+\pi_w)\tth(S_{uw})\ge0$. Let $\zeta\in Z$ be a vertex of degree
$\Delta$. By (B1) it has $\Delta-\pi_\zeta\ge1$ neighbours $w$ in $Z$, and $d_w\le\Delta$ gives
$2/S_{\zeta w}\ge1/\Delta$. Hence $\sum_{ZZ}2/S_{uw}\ge(\Delta-\pi_\zeta)/\Delta$. If
$\pi_\zeta\le1$ this exceeds $1-2/\Delta$. If $\pi_\zeta=2$ it is at least $1-2/\Delta$ and
$E\ge2\,\tth(S_{\zeta w})>0$ for a neighbour $w\in Z$ of $\zeta$. In both cases
\eqref{eq:final} gives $\Sigma_\Gamma>0$.
\end{proof}

\begin{proof}[Proof of Proposition~\ref{prop:block} when $\Delta=2$]
Now every vertex of $Z$ has degree two. We first note:

\smallskip\noindent
($\ast$) \emph{if $a\in\Pi$ and $u\in N(a)\cap Z$, then $a$ and $u$ have no common neighbour in
$Z$.} Indeed, if $w\in Z$ were adjacent to both, then $N(u)=\{a,w\}$ and $N(w)=\{a,u\}$, so
$\{u,w\}$ would be a component of $\Gamma-a$ not containing the other pole, contradicting the
$2$-connectedness of $\Gamma$.

\smallskip
By ($\ast$) and the last statement of Lemma~\ref{lem:deletion}, for $a\in\Pi$ and
$u\in N(a)\cap Z$,
\begin{equation}\label{eq:exact}
 \omega_Z(au)=(x_a+\varepsilon+\pi_u)\,\tht(d_a+2).
\end{equation}
We show that every bracket
$\Lambda(a,u):=\omega_Z(au)-O_a-\varepsilon\tth(d_{a_1}+d_{a_2})$ in Lemma~\ref{lem:lower} is
nonnegative; note $\tth(d_{a_1}+d_{a_2})\le\tth(d_a+2)$ because $d_{a'}\ge2$.

\emph{Poles with $O_a\le x_a\tht(d_a+2)$} (that is, $x_a=0$ or all outside neighbours of degree
at least two). By Lemma~\ref{lem:lower},
$\Lambda(a,u)\ge(x_a+\varepsilon+\pi_u-1)\tth(d_a+2)-x_a\tht(d_a+2)-\varepsilon\tth(d_a+2)
=(\pi_u-1)\tth(d_a+2)+x_a\bigl[\tth(d_a+2)-\tht(d_a+2)\bigr]\ge0$, with strict inequality if
$x_a\ge1$.

\emph{Poles with a single outside neighbour of degree one.} Here $x_a=1$,
$O_a=\tht(d_a+1)$, $d_a=n_a+1+\varepsilon\ge3$, and by \eqref{eq:exact} and $\pi_u\ge1$,
\[
 \Lambda(a,u)\ \ge\ (2+\varepsilon)\tht(d_a+2)-\tht(d_a+1)-\varepsilon\,\tth(d_a+2)
 =\frac{2(d_a-2-\varepsilon)}{d_a(d_a+1)(d_a+2)}=\frac{2(n_a-1)}{d_a(d_a+1)(d_a+2)}\ \ge0 ,
\]
because $n_a+\varepsilon\ge2$; the inequality $\Lambda(a,u)>0$ is strict unless
$\varepsilon=1$, $n_a=1$ and $\pi_u=1$.

Hence $\Sigma_\Gamma\ge\sum_{ZZ}2/S_{uw}+\sum_{a}\sum_u\Lambda(a,u)\ge0$, and it remains to
find one strictly positive term. If some pole has $x_a\ge1$ and all its outside neighbours
have degree at least two, the corresponding $\Lambda(a,u)$ are positive. If some pole $a$ has
an outside neighbour of degree one, then $\Lambda(a,u)>0$ unless $\varepsilon=1$, $n_a=1$ and
$\pi_u=1$ for the unique $u\in N(a)\cap Z$; in that case the second neighbour of $u$ lies in
$Z$, so there is a $ZZ$-edge, of weight $2/4$. Finally, if $x_{a_1}=x_{a_2}=0$, then by
\eqref{eq:exact} $\Lambda(a,u)=(\varepsilon+\pi_u)\tht(d_a+2)-\varepsilon\tth(d_{a_1}+d_{a_2})
\ge(1+\varepsilon)\tht(d_a+2)-\varepsilon\tth(d_a+2)$, which equals $\tht(d_a+2)>0$ for
$\varepsilon=0$ and $\frac{2(d_a-1)}{d_a(d_a+1)(d_a+2)}>0$ for $\varepsilon=1$. In all cases
$\Sigma_\Gamma>0$.
\end{proof}

\section{Proof of the main results}\label{sec:proof}

\begin{proof}[Proof of Theorem~\ref{thm:main}]
For the path with $n\ge3$ vertices, $\Hm(P_n)=2\cdot\frac23+(n-3)\cdot\frac12=\frac{n-1}2+\frac13$,
and $\Hm(K_2)=1>\frac12+\frac13$; so the theorem holds for paths. We prove, by induction on
$n\ge2$, that every connected graph $G$ of order $n$ which is not a path satisfies
$\Hm(G)>\frac12\D(G)+\frac13$. For $n=2$ there is nothing to prove. Let $n\ge3$, let $G$ be
connected and not a path, and assume that every connected graph $G'$ with at least two and
fewer than $n$ vertices satisfies $\Hm(G')\ge\frac12\D(G')+\frac13$ (for paths this is the
computation above, for the other graphs it is the induction hypothesis). Fix $p,q$ with
$d(p,q)=\D=\D(G)$ and a shortest path $P$ from $p$ to $q$.

\emph{Case 1: $G$ has a hanging part,} that is, \eqref{eq:H} fails. Then there are a vertex $v$
and a component $X$ of $G-v$ with $X\cap V(P)=\emptyset$. As $P$ has a vertex other than $v$,
and $G-v$ has a component containing it, $v$ has a neighbour outside $X$. By
Lemma~\ref{lem:hanging}, $\Hm(G)>\Hm(G-X)$. The graph $G-X$ is connected, has between $2$ and
$n-1$ vertices and contains $P$; distances in a subgraph are at least the distances in $G$, so
$\D(G-X)\ge d_{G-X}(p,q)\ge\D$. By the induction hypothesis,
$\Hm(G)>\Hm(G-X)\ge\frac12\D(G-X)+\frac13\ge\frac12\D+\frac13$.

\emph{Case 2: \eqref{eq:H} holds.} If every block of $G$ is a $K_2$, then $G$ is a tree, and
in a tree every vertex of degree at least two is a cut vertex; by Lemma~\ref{lem:chain}(i)
every such vertex is an internal vertex of $P$ and every leaf is $p$ or $q$, so $V(G)=V(P)$
and $G=P$ is a path, contrary to the assumption.
Hence $G$ has a block $\Gamma$ with at least three vertices. Let $\Pi$ be its set of poles and
$Z=V(\Gamma)\setminus\Pi$. By Lemma~\ref{lem:chain}, the hypotheses (B1), (B2) of
Section~\ref{sec:block} hold, so Proposition~\ref{prop:block} yields $z\in Z$ with
$\Hm(G-z)<\Hm(G)$. By Lemma~\ref{lem:chain}(iii), $z$ is not a cut vertex and
$z\notin\{p,q\}$. Thus $G-z$ is connected with $n-1\ge2$ vertices and contains $p$ and $q$,
so $\D(G-z)\ge d_{G-z}(p,q)\ge\D$, and the induction hypothesis gives
$\Hm(G)>\Hm(G-z)\ge\frac12\D(G-z)+\frac13\ge\frac12\D+\frac13$.

This completes the induction. Together with the values for paths, it proves the inequality
and shows that equality holds exactly for the paths with at least three vertices.
\end{proof}

\begin{proof}[Proof of Corollary~\ref{cor:liu}]
Let $n\ge3$ and $\D=\D(G)\le n-1$. By Theorem~\ref{thm:main},
\[
 \Hm(G)-\Bigl(\D+\frac56-\frac n2\Bigr)\ \ge\ \frac{n-1-\D}{2}\ \ge0,\qquad
 \Hm(G)-\Bigl(\frac12+\frac1{3(n-1)}\Bigr)\D\ \ge\ \frac13\Bigl(1-\frac{\D}{n-1}\Bigr)\ \ge0 .
\]
Equality in either line forces $\D=n-1$, i.e.\ $G\cong P_n$, and then equality holds by the
computation of $\Hm(P_n)$.
\end{proof}

\begin{remark}[why averaging]\label{rem:rules}
It is tempting to prove Theorem~\ref{thm:main} by deleting a canonically chosen vertex off the
path $P$, for instance one of minimum degree among the vertices not on $P$. For all connected
graphs of order at most $9$ this vertex does satisfy $\Hm(G-z)<\Hm(G)$, but the rule fails for
large graphs. Take a path $v_0\cdots v_8$, a complete graph $K_{1000}$ and a vertex $z$
adjacent to $v_4$ and to seven vertices of the complete graph. The path is diametral, $z$ has
the smallest degree among the vertices off the path, and $\Hm(G)-\Hm(G-z)\approx-0.0078<0$:
deleting $z$ raises the weights of the two path edges at $v_4$ by
$2(\frac12-\frac25)=\frac15$, while the edge $zv_4$ had weight only $\frac2{11}$ and the
other edges at $z$ are negligible. Similarly, the sum of $\Hm(G)-\Hm(G-z)$ over \emph{all}
vertices that are neither cut vertices nor $p,q$ can be negative: for the graph obtained from
the cycle $C_{400}$ by attaching a pendant vertex to every cycle vertex and adding six
vertices adjacent to all cycle vertices, this sum is about $-0.71$. The two reductions of the
proof are designed to avoid these phenomena: hanging parts are removed first, and inside a
block the losses at the two poles are controlled by Lemma~\ref{lem:pole}.
\end{remark}

\begin{remark}[sharpness]
Equality holds only for paths. Among the connected graphs of order at most $11$ that are not
paths, the minimum of $\Hm(G)-\frac12\D(G)-\frac13$ equals $\frac7{30}$ for every order
$5\le n\le11$; it is attained by the path with one pendant vertex attached to the neighbour of
an end vertex. For $n=4$ the minimum is $\frac16$ (the star $K_{1,3}$).
\end{remark}

\section{Computational checks}\label{sec:computation}

The proof above does not rely on machine computation; the only numerical comparisons in it
(Table~\ref{tab:pole} and the small cases in Lemma~\ref{lem:pole}) involve explicit fractions.
Independently, we checked the following statements with exact rational arithmetic (Python
\texttt{fractions}) or, for the largest enumerations, in floating point with a margin far
above rounding errors; graphs were enumerated with \texttt{geng} from
\texttt{nauty}~\cite{nauty}. The programs accompany this manuscript.
\begin{itemize}
\item Theorem~\ref{thm:main} for all connected graphs of order at most $11$
 ($1{,}006{,}700{,}565$ graphs of order $11$): the inequality holds, equality occurs only for
 paths, and the minimum over non-paths is $\frac7{30}$ (order $4$: $\frac16$). The inequality
 also holds for about $8{,}000$ random and structured graphs with up to several hundred
 vertices.
\item Lemma~\ref{lem:deletion} and Lemma~\ref{lem:hanging} on all connected graphs of order at
 most $8$: the formula for every vertex subset arising in the proof, and the strict
 inequality $\Hm(G)>\Hm(G-X)$ for every vertex $v$ and every admissible component $X$.
\item The chain of inequalities of Section~\ref{sec:block} for every block with at most two
 cut vertices of every connected graph of order at most $8$ and every choice of two poles
 containing these cut vertices ($219{,}088$ cases at order $8$, of which $217{,}539$ satisfy
 \textup{(B2)}): $\Sigma_\Gamma$ is at least the bound of Lemma~\ref{lem:lower}, which is at
 least $\sum_{ZZ}2/S-\Bad_{a_1}-\Bad_{a_2}$; under \textup{(B2)},
 $\Bad_a\le\frac12-\frac1\Delta$, and the smallest value of $\Sigma_\Gamma$ is
 $0.2769\ldots$
\item Lemma~\ref{lem:pole} for all $n,x\le200$ and a range of $\Delta$ up to $10^5$; equality
 occurs only in the case stated.
\item The induction step itself: for every connected non-path graph of order at most $8$,
 every diametral pair and every shortest path between its two vertices ($147{,}561$ instances
 at order $8$), either a hanging part exists or the structure of Lemma~\ref{lem:chain} holds
 and each block with at least three vertices contains a vertex $z$ as in the proof.
\item The two examples of Remark~\ref{rem:rules}.
\end{itemize}

\section{Status of the result}\label{sec:status}

\paragraph{What is claimed.} Theorem~\ref{thm:main} and Corollary~\ref{cor:liu} are claimed
as theorems, with complete proofs in Sections~\ref{sec:deletion}--\ref{sec:proof}. The only
external result used is Zhong's inequality \eqref{eq:zhong} (in Lemma~\ref{lem:hanging}),
together with standard facts about blocks. The manuscript has not been refereed.

\paragraph{The criteria.}
\begin{center}
\begin{tabular}{@{}lp{0.73\textwidth}@{}}
\toprule
Criterion & Assessment\\
\midrule
G1 named problem & Liu's conjecture on the harmonic index and the diameter \cite{Liu} (2013), restated in \cite{JM1,JM2,AI}.\\
G2 original statement & Corollary~\ref{cor:liu} is the conjecture as stated, for all $n\ge3$, with the same constants; nothing is transformed.\\
G3 open & To our knowledge the conjecture was proved only for trees \cite{Liu}, unicyclic graphs \cite{JM1,JM2} and quasi-tree graphs \cite{AI}. Our search (web search engines, arXiv, and the abstracts and introductions of the papers cited here) found no proof for general graphs; a bounded search cannot exclude one.\\
G4 complete proof & The proof is unconditional and covers every connected graph. Computations are used only as an independent check.\\
G5 non-restatement & The result is the new, order-free inequality $\Hm\ge\frac12\D+\frac13$; the proof uses new estimates (Lemmas~\ref{lem:hanging} and~\ref{lem:pole}, Proposition~\ref{prop:block}), not a change of variables.\\
G6 priority & We found neither Theorem~\ref{thm:main} nor the deletion formula of Lemma~\ref{lem:deletion} in the literature, with the same caveat as for G3.\\
\bottomrule
\end{tabular}
\end{center}

\paragraph{Relation to earlier work in this series.} The manuscript \cite{VMRC} proves that
$\Hm(T)$ is at least the average distance of $T$ for every tree and reduces the general
statement, a strengthening of Fajtlowicz's conjecture, to a single inequality which remains
open. The present paper uses the same philosophy (exact first-order identities for a
degree-based index under deletions) but for vertices instead of edges, and here the method
closes completely.

\paragraph{Questions.} (1) Is there a constant $c>\frac12$ and a natural class of graphs (for
instance $2$-connected graphs) with $\Hm(G)\ge c\,\D(G)$? (2) Jerline and Michaelraj
\cite{JM1,JM2} and Abdolghafourian and Iranmanesh \cite{AI} obtained the sharper bound
$\Hm(G)\ge\D(G)+\frac53-\frac n2$ for unicyclic and quasi-tree graphs (with a few small
exceptions); determine the minimum
of $\Hm(G)-\frac12\D(G)$ over all connected graphs that are not trees. (3) Does the
vertex-deletion calculus settle the analogous conjectures relating the harmonic index to the
radius?

\begin{thebibliography}{99}
\bibitem{AI} A.~Abdolghafourian and M.~A.~Iranmanesh, On harmonic index and diameter of
quasi-tree graphs, \emph{J. Math.} \textbf{2021}, Article ID 6650407; arXiv:2011.05752.
\bibitem{AZG} A.~Ali, L.~Zhong and I.~Gutman, Harmonic index and its generalizations: extremal
results and bounds, \emph{MATCH Commun. Math. Comput. Chem.} \textbf{81} (2019), 249--311.
\bibitem{Faj} S.~Fajtlowicz, On conjectures of Graffiti II, \emph{Congr. Numer.} \textbf{60}
(1987), 187--197.
\bibitem{FMS} O.~Favaron, M.~Mah\'eo and J.-F.~Sacl\'e, Some eigenvalue properties in graphs
(conjectures of Graffiti II), \emph{Discrete Math.} \textbf{111} (1993), 197--220.
\bibitem{JM1} J.~Amalorpava Jerline and L.~Benedict Michaelraj, On a conjecture of harmonic
index and diameter of graphs, \emph{Kragujevac J. Math.} \textbf{40} (2016), 73--78.
\bibitem{JM2} J.~Amalorpava Jerline and L.~Benedict Michaelraj, On harmonic index and diameter
of unicyclic graphs, \emph{Iranian J. Math. Sci. Inform.} \textbf{11} (2016), 115--122.
\bibitem{Liu} J.~Liu, On harmonic index and diameter of graphs, \emph{J. Appl. Math. Phys.}
\textbf{1} (2013), no.~3, 5--6.
\bibitem{nauty} B.~D.~McKay and A.~Piperno, Practical graph isomorphism, II, \emph{J. Symbolic
Comput.} \textbf{60} (2014), 94--112.
\bibitem{VMRC} Verifiable Mathematical Reasoning Corpus, The harmonic index of a tree
dominates its average distance, research manuscript, 15 November 2026.
\bibitem{Zhong} L.~Zhong, The harmonic index for graphs, \emph{Appl. Math. Lett.} \textbf{25}
(2012), 561--566.
\end{thebibliography}

\end{document}
